\section{Introduction}\label{sec:introduction}

Ideal classes in a totally real number field give compact orbits of a
diagonal group on a space of lattices.  A natural arithmetic question is
whether the finite collections of orbits obtained in this way become
uniformly distributed as their discriminants grow.  We prove this for the
complete packets attached to invertible ideals of every order in a primitive
totally real quartic field.  The order may have arbitrarily large index in
the ring of integers.

Let
\[
 X_4=\mathrm{SL}_4(\Z)\backslash\mathrm{SL}_4(\R)
\]
be the space of covolume-one row lattices in $\R^4$, and let $m_4$ be its
invariant probability measure.  Put
\[
 H_4=\{u\in\R^4:u_1+u_2+u_3+u_4=0\},\qquad
 A_4=\{\diag(e^{u_1},e^{u_2},e^{u_3},e^{u_4}):u\in H_4\}.
\]
The group $A_4$ acts on $X_4$ by right multiplication.  Its Haar measure
is $du_1\,du_2\,du_3$, with $u_4=-u_1-u_2-u_3$.

A quartic field $K$ is \emph{primitive} if it has no field strictly between
$\Q$ and $K$.  Let $K$ be primitive and totally real, let $\OO\subseteq\OO_K$
be an order, and fix an ordering $\sigma=(\sigma_1,\ldots,\sigma_4)$ of its
real embeddings.  Here an order is a full-rank subring containing $1$.
A fractional $\OO$-ideal $I$ is \emph{proper} if its multiplier order
$\{a\in K:aI\subseteq I\}$ is $\OO$, and it is \emph{invertible} if
$IJ=\OO$ for some fractional $\OO$-ideal $J$.  The ordinary ideal class
group $\Pic(\OO)$ consists of proper invertible fractional $\OO$-ideals
modulo $I\sim aI$ for $a\in K^\times$.

For such an ideal put
\begin{equation}\label{eq:intro-lattice}
 \Lambda_{I,\sigma}
   =\covol(\sigma(I))^{-1/4}\sigma(I)\in X_4.
\end{equation}
A covolume-one lattice is represented here by any positively oriented
basis; two such bases differ by an element of $\mathrm{SL}_4(\Z)$.
Let $W_{\mathrm{all}}$ be the sixteen diagonal matrices with entries in
$\{1,-1\}$.  Choose one ideal in each ordinary class and form the finite
set of distinct orbits
\[
 \mathcal P_{\OO,\sigma}
 =\{\Lambda_{I,\sigma}wA_4:I\text{ a class representative},\
                                      w\in W_{\mathrm{all}}\}.
\]
A matrix in $W_{\mathrm{all}}$ acts on the underlying lattice, so a sign
matrix of determinant $-1$ also gives a point of $X_4$.  Using all signs
makes the set independent of the chosen ordinary ideal representatives.
Every orbit is compact by the unit theorem; Section~\ref{sec:packets}
proves these assertions and describes the associated measure.

For an orbit $\mathcal Y\in\mathcal P_{\OO,\sigma}$, let $\nu_{\mathcal Y}$
be its invariant probability and let $\vol(\mathcal Y)$ be its volume for
the specified Haar measure.  Define
\begin{equation}\label{eq:intro-measure}
 \mu_{\OO,\sigma}
 =\frac{\displaystyle\sum_{\mathcal Y\in\mathcal P_{\OO,\sigma}}
              \vol(\mathcal Y)\nu_{\mathcal Y}}
        {\displaystyle\sum_{\mathcal Y\in\mathcal P_{\OO,\sigma}}
              \vol(\mathcal Y)}.
\end{equation}
Each distinct orbit occurs once in these sums.

\begin{theorem}\label{thm:main}
Let $(K_i,\OO_i,\sigma_i)$ be any sequence in which $K_i$ is a primitive
totally real quartic field, $\OO_i$ is an order in $K_i$, and $\sigma_i$ is
an ordering of its real embeddings.  If $|\Disc(\OO_i)|\to\infty$, then
\[
 \int_{X_4} f\,d\mu_{\OO_i,\sigma_i}
       \longrightarrow \int_{X_4} f\,dm_4
       \qquad\text{for every }f\in C_c(X_4).
\]
There is no escape of mass: for every $\varepsilon>0$ there is a compact
$\mathcal K\subset X_4$ such that
$\mu_{\OO_i,\sigma_i}(\mathcal K)\ge1-\varepsilon$ for all sufficiently
large $i$.
\end{theorem}

In particular, the field may remain fixed while the order index tends to
infinity.  The packet in Theorem~\ref{thm:main} uses every class in
$\Pic(\OO)$, so its ideals are locally free of rank one over the specified
order.  The proof uses this complete class average in both of its counting
arguments.

\subsection{History and context}

For real quadratic fields, compact diagonal orbits are closed geodesics on
the modular surface.  Linnik's ergodic method and Skubenko's theorem for
positive discriminants established arithmetic distribution under the
condition $\left(\frac d p\right)=1$ for the discriminant $d$ at a fixed
odd prime $p$ \cite{Linnik1968,Skubenko1962}.
Using Iwaniec's estimates for half-integral weight forms \cite{Iwaniec1987},
Duke removed that restriction for the geodesic collections associated with
positive fundamental discriminants \cite[Theorem~1]{Duke1988}.
Chelluri extended the corresponding statement to the unit tangent bundle,
as recorded in \cite[Section~1.2, p.~253]{ELMV2012}.
Einsiedler, Lindenstrauss, Michel, and Venkatesh later gave a proof using
entropy for all positive nonsquare order discriminants
\cite[Theorem~2.3]{ELMV2012}.

In higher degree the diagonal group has higher rank, and the arithmetic
packets lie in spaces of higher-dimensional lattices.  Einsiedler,
Lindenstrauss, Michel, and Venkatesh developed the discriminant and volume
framework connecting separation of periodic torus orbits to entropy
\cite{ELMV2009}.  Their cubic theorem proves equidistribution of packets
that include those arising from nonmaximal totally real cubic orders
\cite[Theorem~1.4]{ELMV2011}.  The local homothety class of a lattice is
part of their packet data; a multiplier order alone need not specify it
\cite[Section~5.5]{ELMV2011}.  Their higher-dimensional discussion also
identifies the possibility of proper intermediate homogeneous limits
\cite[Section~1.6.3]{ELMV2011}.  The theorem of Einsiedler, Katok, and
Lindenstrauss makes a diagonal-ergodic probability with positive entropy
homogeneous \cite[Theorem~1.3]{EKL2006}; in prime dimension at least three
its homogeneous conclusion is already Haar \cite[Corollary~1.4]{EKL2006}.

For each fixed prime $p\ge5$, Lemke Oliver, Thorner, and Zaman proved
equidistribution for the orbit collections attached to orders in totally
real degree-$p$ fields as the order discriminant grows, outside an
exceptional set containing $O_{p,\varepsilon}(Q^\varepsilon)$ fields of
field discriminant at most $Q$ \cite[Theorem~2.3]{LOTZ2024}.
The separate preprint \cite[Theorem~1.1]{PrimeDegreePackets} treats every
packet in each fixed prime degree at least five and permits arbitrary local module types.
We adapt its ordinary-vector method in the argument below.
For quartic packets, Khayutin obtained entropy bounds for probability
weak-$*$ limits of maximal-order packet measures with two-transitive Galois
groups, including the primitive quartic types $\mathfrak A_4$ and $S_4$
\cite[Theorem~1.1]{Khayutin2019Tori}.  Wieser and Yang obtain entropy bounds
for certain maximal-order quartic packets whose fields have a quadratic
subfield, under further hypotheses \cite{WieserYang2026}.  The present
argument must establish tightness and positive entropy in almost every
diagonal-ergodic component for arbitrary orders, and then exclude the
proper homogeneous limits allowed in degree four.

The arithmetic use of a sieve to rule out intermediate algebraic limits
has a precedent in Khayutin's work on joint distributions of CM points
\cite{Khayutin2019CM}.  Local Fourier analysis for nonmaximal orders already
appears in the cubic work of Einsiedler, Lindenstrauss, Michel, and Venkatesh
\cite[Proposition~9.9]{ELMV2011}.  Here the object being transformed is a
specific quadratic product map on an exterior lattice.  The resulting
estimate controls arbitrary conductor depth in a cubic-resolvent count,
before the remaining multiplicative weights are inserted by a sieve.

\subsection{The two counts and the conductor estimate}

Write
\[
 \begin{gathered}
 R=\OO_K,\qquad D=|\Disc(K)|,\qquad q=[R:\OO],\\
 X=q\sqrt D=\sqrt{|\Disc(\OO)|},\qquad
 \kappa=\Res_{s=1}\zeta_K(s).
 \end{gathered}
\]
We also write $a_K(n)$ for the number of nonzero integral $R$-ideals of
norm $n$.  Both counting arguments preserve the factor $\sqrt D\,\kappa$
that normalizes the packet after its finite local labels are averaged.

Sections~\ref{sec:packets} and~\ref{sec:ordinary} treat ordinary lattice
vectors.  An invertible $\OO$-ideal is locally a unit translate of $\OO$
inside a maximal-order ideal.  Averaging those translates gives an exact
unfolding into a weighted sum over integral $R$-ideals.  A local moment
bound controls the order index, while Stark's exceptional-zero result and
the absence of a quadratic subfield retain the residue $\kappa$ in the
ideal count.  Counting short vectors proves tightness.  Counting vectors
in boxes away from the coordinate hyperplanes gives an $O(r^4)$ bound for
radius-$r$ neighborhoods in $X_4$, uniformly on each compact set of centers
in every probability limit.  For a diagonal flow with distinct weights,
a two-sided tube is covered by $O(r^{-3})$ such neighborhoods with centers
in a fixed compact enlargement, where $r$ decreases exponentially with the
tube length.  The resulting exponential decay also holds for normalized
averages over sets of $A_4$-ergodic components with positive measure,
because they are dominated by a multiple of the limit measure.  This forces
positive entropy in almost every component.

Section~\ref{sec:rigidity} applies measure classification and identifies
the remaining obstruction.  A proper positive-entropy component has two
coordinate blocks of dimension two.  Each lattice in its supporting orbit
then has a vector in its exterior square supported on the two opposite
coordinates $ij$ and $kl$, with both coefficients nonzero.  There are three
choices of the partition $ij|kl$.

For an exterior vector $v=(v_{ij})_{i<j}$, the positive diagonal action
preserves the three signed products
\[
 (v_{12}v_{34},\,-v_{13}v_{24},\,v_{14}v_{23}).
\]
Their sum is an integer on the exterior square of a unimodular lattice.
For the arithmetic lattices in our packets, Section~\ref{sec:exterior}
places the products in the three embeddings of a cubic resolvent field
$F$.  Fixing their sum at a nonzero integer $m$ puts them on a trace-$m$
plane.  If the four exterior coordinates outside $ij,kl$ have size at most
$r$, two product coordinates have size at most $r^2$.  Thus the geometric
obstruction is tested by a rectangle of area $O(r^4)$ on that plane.
Theorem~\ref{thm:exterior-count} supplies the corresponding arithmetic
rectangle bound, with the multiplicity of every product and the packet
normalization included.

The main new work lies in proving that bound uniformly in $\OO$.  At primes
dividing $Dq$, the difficult local factors in the multiplicity are given by
the density of the pushforward of Haar measure on the exterior square of the
order lattice, with a unit scalar, under the quadratic product map.
Its exact fiber measure accounts for ramification and for the order index.
An element of $F$ acts on each opposite-coordinate plane by the value of the
corresponding real embedding; the resulting action descends to the rational
exterior space $\bigwedge^2_{\Q}K$.  The elements preserving
$\bigwedge^2_{\Z}\OO$ form an order $C_F$.  Its discriminant is bounded
below by a fixed positive power of $X$ and above by a fixed power of $X$,
even if the field $F$ remains fixed.

Section~\ref{sec:local-fourier} analyzes these densities through their
quadratic Fourier transforms.  Denominator ideals organize the frequencies;
the stabilizer order controls the decay of the Fourier coefficients and the
support and size of the corresponding inverse transforms.  The local proof
is expressed in terms of the quadratic lattice, its action by a cubic
\'etale algebra, and its stabilizer order, a formulation suitable for
analogous cubic-resolvent counts.  In Section~\ref{sec:additive}, Poisson
summation gives the trace density from scalar frequencies.  Complementary
estimates control the total nonscalar contribution using discriminant
separation, bounds for the Fourier and inverse transforms, and congruences
supplied by inverse quadratic forms.  The estimate is uniform for small
residue moduli and for shrinking rectangles.

At the remaining primes, elementary valuation records give a local factor
with mean $1+(a_K(p)-1)/p+O(p^{-2})$.  Section~\ref{sec:sieve} averages
these factors by a weighted Selberg sieve.  The subtraction of $1$ in the
local mean removes the logarithm in the prime-sum bound, leaving $\kappa$.
Finally, Section~\ref{sec:conclusion} unfolds exterior vectors along the
diagonal group.  The fiber lengths are logarithmic and summable over the
rectangles above.  The resulting $O(r^4)$ bound excludes every proper
homogeneous component and proves Theorem~\ref{thm:main}.

\subsection{Conventions}

The notation $\Nm_{E/k}$ denotes the algebraic field norm.  Ideal norms
$\Nm(\mathfrak a)$ are positive absolute norms, extended multiplicatively
to fractional ideals; absolute values are written explicitly for field
norms when needed.  For a rational prime $p$,
a subscript $p$ denotes tensoring a field with $\Q_p$ or a lattice with
$\Z_p$.  For a prime $\mathfrak p$ or $\mathfrak q$ of a local field we use
the integer valuation with value $1$ on a uniformizer, and write $e$ and $f$
for its ramification index and residue degree.  A normalized extension of
$v_p$ has $v_p(p)=1$, and $|p|_p=p^{-1}$.  Ideal norms include the residue
degrees.  We use $\mathfrak p$ for primes of $K$ and $\mathfrak q$ for primes
of $F$ when the two occur together.

All logarithms are natural, and $\log_+t=\max(0,\log t)$ for $t>0$.
An upper bound $X^{o(1)}$ means $O_\eta(X^\eta)$ for every fixed $\eta>0$.
Unless stated otherwise, constants are uniform in $K$, $\OO$, and the
ordering of the embeddings.  In the exterior estimates a nonzero trace $m$
and a bounded real region are fixed before $X$ tends to infinity; permitted
dependence on them is stated in the relevant results.
